\documentclass[11pt]{article}
\usepackage{lmodern}          % AER font
\usepackage[T1]{fontenc}      % Better font encoding
\usepackage{amsmath, amssymb}
\usepackage[margin=1in]{geometry}
\usepackage{graphicx}
\usepackage{url}
\usepackage[round]{natbib}
\usepackage{bibentry}         % lets \nobibliography resolve cites without printing a list
\usepackage{hyperref}
\usepackage{fancyhdr}

%=======================================
% Format helpers
%=======================================
% Target school -- used in the header and the body so they stay in sync.
\newcommand{\school}{[University name]}
\newcommand{\program}{[Program name]}

% Run-in section heading: bold label ending in a period, body text follows
% on the same line (matches the reference format).
\newcommand{\shead}[1]{\par\vspace{0.85em}\noindent\textbf{#1.}\ }

% Indented paragraphs, no extra vertical space between them.
\setlength{\parindent}{1.5em}
\setlength{\parskip}{0pt}

% Footer: page number centred, word count on the right. The count is read
% from Rhea_PS.wc, which the "texcount -> pdflatex" recipe regenerates on each
% build. Text between %TC:ignore and %TC:endignore is excluded from the count.
\newcommand{\wordcount}{\IfFileExists{\jobname.wc}{\input{\jobname.wc}\unskip}{??}}
\pagestyle{fancy}
\fancyhf{}
\renewcommand{\headrulewidth}{0pt}
\fancyfoot[C]{\thepage}
\fancyfoot[R]{\small Word count: \wordcount}


\begin{document}

%TC:ignore
%=======================================
% Header
%=======================================
\noindent
\begin{minipage}[t]{0.60\textwidth}
  {\large\bfseries Personal Statement}
\end{minipage}\hfill
\begin{minipage}[t]{0.38\textwidth}
  \raggedleft
  Rhea Choudhary\\[2pt]
  \school
\end{minipage}

\vspace{0.5em}
\noindent\rule{\textwidth}{0.8pt}

\vspace{1.1em}
%TC:endignore

%=======================================
% Body
%=======================================

%my sop research agenda: how large-scale government programs can better achieve their desired policy objectives, and the externalities created by such programs in the pursuit of these goals

%prompts for personal statement:
 %This statement should be a discussion of the journey that has led to your decision to seek a graduate degree.

 %Pivotal experiences, opportunities, and/or challenges that have influenced your educational and professional development.

%Describe a personal experience that influenced your decision to pursue graduate study. Explain how the lessons from this experience would enrich Princeton's residential scholarly community.

When I was about eight years old, my parents decided to move from the Middle East to Australia to secure a better and brighter future for me and my brother. However, my dad struggled to find work once we got there, so my parents made the tough decision to live separately -- my mum with my brother and I in Australia, and my dad overseas. For almost two decades, my mum effectively functioned as a single parent and I watched as she pulled every string at her disposal to give us the future they so desperately wanted for us. Back in the Middle East, my dad worked long hours to send us money and afford a few trips a year to visit, including a special trip for my 16th birthday, an occassion he usually had to miss.
\\~\\
As I look back on my high school and university experiences, I think a lot about how important my family and the informal safety net they provided was to my academic and professional journey. I could take risks in my career, take on more study or choose more demanding professional roles in large part because I had informal support. But I also think a lot about how getting through the door to these opportunities required government support; I went to a public university, used government-issued university loans, and relied income assistance when studying full-time. In this way, both government assistance and family support played complementary roles in helping my career. But I've also seen examples where both fall short. For instance, I watched as my brother was unemployed for over three years and needed more than just the imperfect income assistance provided by unemployment insurance. Informal support can only go so far -- moral support and encouragement from my parents can often be viewed as patronising or insensitive, and financial support to wedge the gap that unemployment insurance creates can only last so long. Like most young Australians, he is college educated, so mutual obligation requirements that encourage short-term courses didn't work for him. Job-search in fields like IT and computer science, which is his specialisation, is incredibily challenging because the nature of these jobs are rapidly changing to due AI. Government programs aren't well-positioned to alleviate labour-market scarring in the face of such technological disruption. 
\\~\\
These experiences have strongly motivated my decision to pursue graduate study, because I want to answer important questions like these that are front-of-mind within the academic community as well as in policymaking. I want to grow my knowledge set on key developments and techniques in the literature, and eventually push the frontier of the discipline to answer questions of deep importance to me.  

%These experiences have motivated my research interests and agenda, which lie in empirically researching large-scale government programs and its broader welfare effects. In my graduate studies I want to extend this agenda to consider both spillovers of large-scale government programs, when or how informal safety nets act as substitutes or complements for formal support, and how this interacts with program objectives. For example, I want to explore why some formal assistance programs crowd out existing informal support with no net welfare gain \citep[e.g.,][]{havnes_mogstad2015}, while others may create complementarities within the household \citep{rennane2020}.







%.
%Something about:
%- the informal social safety net -> the informal safety net that my parents provided had important spillover effects on their own life -- effectively that their own lives had to be put on hold.
%How important the informal safety net is in informing individual decisions, and how xxx

%What unique worldview could I bring?
%from my personal experience: 
%- 

%from my professional experience:
%- how policymaking is made, how policymakers trade off costs and benefits and what drives their decisions. This has given me a unique perspective on how to frame research questions and how to communicate findings in a way that is useful for policymakers, and also how to pitch research that require their data or support.

%=======================================

%TC:ignore
% Citations resolve against bib_PS.bib, but no reference list is printed.
\bibliographystyle{abbrvnat}
\nobibliography{bib_PS}
%TC:endignore

\end{document}
